% !PRINTED 362
\weberpage{362}

Setzt man also
\weq{5}{a_1 = N(\mathfrak{m}_1) = \mathfrak{m}_1\mathfrak{m}_1',\quad \frac{b_1+\sqrt{D}}{2} \equiv 0\ (\mathrm{mod}\ \mathfrak{m}_1),}
so erhält man eine andere quadratische Form $(a_1,b_1,c_1)$, von der wir nachweisen können, daß sie mit $(a,b,c)$ äquivalent ist.

Wir setzen zur Abkürzung:
\[
\omega = \frac{b+\sqrt{D}}{2a},\quad \omega_1 = \frac{b_1+\sqrt{D}}{2a_1}
\]
und bemerken, daß
\[
\eta a_1 = \eta \mathfrak{m}_1\mathfrak{m}_1' = \mathfrak{m}\mathfrak{m}_1',
\]
\[
\eta \omega_1 a_1 = \frac{\mathfrak{m}}{\mathfrak{m}_1}\,\frac{b_1+\sqrt{D}}{2}
\]
durch $\mathfrak{m}$ teilbare Zahlen in $[Q]$ sind [nach (5)]. Demnach lassen sich nach \S\,97, 7. die ganzen rationalen Zahlen $\alpha$, $\beta$, $\gamma$, $\delta$ so bestimmen, daß:
\weq{6}{\begin{aligned}\eta a_1&=a(\alpha-\gamma\omega),\\ \eta a_1\omega_1&=a(-\beta+\delta\omega),\end{aligned}}
und man beweist ganz wie in \S\,95, daß
\[
\alpha\delta - \beta\gamma = 1
\]
sein muß, daß also die beiden Formen
\[
(a,b,c),\quad (a_1,b_1,c_1)
\]
äquivalent sind. Wir haben also:

11. \gesperrt{Satz. Jede durch ein Ideal $\mathfrak{a}$ repräsentierte Idealklasse nach $[Q]$ entspricht einer durch die Form $(a,b,c)$ repräsentierten Formklasse $A$ und umgekehrt.}

Um die Formklasse $A$ zu erhalten, nehme man ein primäres Ideal $\mathfrak{m}$ der zu $A$ reziproken Klasse $A^{-1}$ und setze $a = N(\mathfrak{m})$,
$\frac{b+\sqrt{D}}{2} \equiv 0\ (\mathrm{mod}\ \mathfrak{m})$. Und ist umgekehrt die Form $(a,b,c)$ gegeben, so ist der größte gemeinschaftliche Teiler $\mathfrak{m}$ von
\[
a,\quad \frac{b+\sqrt{D}}{2}
\]
ein Ideal der Klasse $A^{-1}$.

\begin{center}
\S\,100. \gesperrt{Idealklassen nach den Ordnungen.}
\end{center}

Der Begriff der Äquivalenz der Ideale und die Sätze über Klassenzahlen lassen sich am übersichtlichsten darstellen, wenn
%CONT>>
